\section{Discrete dynamical systems}
For $x_{k+1}=Ax_k$, start by expanding $x_0$ in an eigenbasis:
$x_0=\sum_j c_jv_j$, $Av_j=\lambda_jv_j$. Then
\[x_k=\sum_j c_j\lambda_j^kv_j.\]
This follows by induction: multiplication by $A$ adds one power of
each eigenvalue. A mode with $|\lambda|<1$ decays; one with
$|\lambda|>1$ grows. A negative eigenvalue also reverses its
mode's sign at each step. The direction of fastest decay corresponds
to the smallest modulus, whereas fastest growth corresponds to the
largest modulus. These discrete-time tests concern modulus, not merely
the sign of an eigenvalue.

\practice{5.6.P1}{Determine the three initial coefficients}
The given eigenvectors are $v_1=(-2,2,1)^T$, $v_2=(2,1,2)^T$,
$v_3=(1,2,-2)^T$. Solve $c_1v_1+c_2v_2+c_3v_3=(1,11,-2)^T$:
\[-2c_1+2c_2+c_3=1,\quad2c_1+c_2+2c_3=11,
\quad c_1+2c_2-2c_3=-2.\]
Elimination gives $(c_1,c_2,c_3)=(2,1,3)$; substitution checks
all three coordinates. Therefore
$x_k=2v_1+(2/3)^kv_2+3(1/3)^kv_3$.

\practice{5.6.P2}{Retain only the nondecaying mode}
The powers $(2/3)^k$ and $(1/3)^k$ tend to zero. Thus the preceding
solution tends to $2v_1=(-4,4,2)^T$, rather than zero. The eigenvalue
$1$ preserves its component indefinitely.

\exercise{5.6.1}{Use the given eigenvectors without reconstructing the matrix}
Solve $c_1(1,1)^T+c_2(-1,1)^T=(9,1)^T$.
The coordinate equations $c_1-c_2=9$, $c_1+c_2=1$ give
$c_1=5,c_2=-4$.
\step{(a) Apply $A$ once to the decomposition.}
$x_1=5(3)(1,1)^T-4(1/3)(-1,1)^T=(49/3,41/3)^T$.
\step{(b) Apply the same rule $k$ times.}
\[x_k=5\,3^k(1,1)^T-4\,(1/3)^k(-1,1)^T.\]
At $k=0$ this gives the specified initial vector; multiplication by
$A$ verifies the recurrence at every subsequent step.

\exercise{5.6.2}{Separate the growing mode from two decaying modes}
The coefficient system $[v_1\ v_2\ v_3]c=x_0$ is
\[c_1+2c_2-3c_3=-2,\quad c_2-3c_3=-5,
\quad-3c_1-5c_2+7c_3=3.\]
It gives $c=\left[\begin{matrix}2\\1\\2\end{matrix}\right]$. Hence
\[x_k=2\cdot3^kv_1+(4/5)^kv_2+2(3/5)^kv_3.\]
The last two modes decay, whereas the first grows. More precisely,
$x_k/3^k\to2v_1=(2,0,-6)^T$. Thus the sequence is unbounded and
its direction approaches the line through $v_1$; its second coordinate
tends to zero while its first and third grow in opposite signs.

\exercise{5.6.3}{Increase predation in the owl--wood-rat model}
For $p=.2$ the matrix is
$A=\begin{bmatrix}.5&.4\\-.2&1.1\end{bmatrix}$.
The characteristic polynomial is $\lambda^2-1.6\lambda+.63$
with roots $.9,.7$. Eigenvectors are $(1,1)^T$ and $(2,1)^T$,
as direct multiplication verifies. Thus
\[x_k=c_1(.9)^k(1,1)^T+c_2(.7)^k(2,1)^T.\]
Both modes decay. Under the exercise's $c_1>0$ assumption, the
long-term positive populations both decline with factor approaching
$.9$ per month. The first component counts owls; the second is the
number of wood rats in thousands, as in Example 1.

\exercise{5.6.4}{The borderline between eventual decline and growth}
With $p=.125$, the matrix has eigenvalues $1,.6$ and eigenvectors
$(4,5)^T,(4,1)^T$. For instance
$A(4,5)^T=(4,5)^T$ and $A(4,1)^T=.6(4,1)^T$.
Consequently
\[x_k=c_1(4,5)^T+c_2(.6)^k(4,1)^T\longrightarrow c_1(4,5)^T.\]
The populations approach constant levels rather than tending to zero
or growing indefinitely. Small parameter changes can move the
largest eigenvalue below or above $1$, changing the prediction to
decline or growth. This explains the sensitivity called an unstable
equilibrium in the exercise; the equilibrium line is not a robust
attracting equilibrium under model changes. Coordinate ratio $4:5$
means 4 owls per 5 thousand rats under the model's units.

\exercise{5.6.5}{Long-term growth and the dominant population ratio}
For $p=.325$, $A=\begin{bmatrix}.4&.3\\-.325&1.2\end{bmatrix}$.
Its eigenvalues solve $\lambda^2-1.6\lambda+.5775=0$, giving
$1.05,.55$. A dominant eigenvector is $(6,13)^T$, because
$A(6,13)^T=(6.3,13.65)^T=1.05(6,13)^T$.
Under the stated positive dominant-coefficient assumption,
$x_k$ is asymptotic to $c_1(1.05)^k(6,13)^T$.
\answer{Both populations eventually grow about $5\%$ per step,
and their coordinate ratio tends to $6:13$. If the prey coordinate
uses the thousands convention of Example 1, this is 6 owls per
13 thousand squirrels; coordinate scaling must be retained when
interpreting a physical ratio.}

\exercise{5.6.6}{Choose the predation rate that preserves population levels}
At $p=.5$, the characteristic polynomial is
$\lambda^2-1.6\lambda+.63$, whose roots are $.9,.7$.
Both powers tend to zero, so both populations perish in this model.
For constant limiting levels, require an eigenvalue $1$:
\[\det(A-I)=(-.6)(.2)-(.3)(-p)=-.12+.3p=0.\]
Thus $p=.4$. The other eigenvalue is $.6$, and a $1$-eigenvector
is $(1,2)^T$. Therefore the solutions tend to $c_1(1,2)^T$.
\answer{The preserving parameter is $p=.4$, and the limiting
coordinate ratio is $1:2$, with prey-unit scaling interpreted as
in the preceding exercise.}

\exercise{5.6.7}{A saddle with one attracting and one repelling direction}
\begin{enumerate}[label=(\alph*)]
\item Since $3>1$ but $1/3<1$, the origin is a saddle point.
\item The line spanned by $v_1=(1,1)^T$ is repelling, with factor
$3$ each step. The line spanned by $v_2=(-1,1)^T$ is attracting,
with factor $1/3$. These are also the greatest repulsion and attraction
directions, respectively.
\item In eigen-coordinates $x=av_1+bv_2$, the update is
$(a,b)\mapsto(3a,b/3)$. The sketch uses those coordinates:
horizontal arrows point away from zero and vertical arrows toward
zero. Off-axis trajectories approach the repelling axis while moving
outward. The curves are schematic connections of discrete points,
not continuous-time trajectories; $ab$ remains constant here.
\end{enumerate}
\begin{center}\setlength{\unitlength}{1pt}
\begin{picture}(220,165)
\put(15,80){\vector(1,0){190}}\put(110,5){\vector(0,1){150}}
\put(197,66){$a$ ($v_1$)}\put(118,147){$b$ ($v_2$)}
\put(110,80){\vector(1,0){60}}\put(110,80){\vector(-1,0){60}}
\put(110,140){\vector(0,-1){35}}\put(110,20){\vector(0,1){35}}
\qbezier(122,150)(124,100)(200,92)
\qbezier(20,92)(96,100)(98,150)
\qbezier(122,10)(124,60)(200,68)
\qbezier(20,68)(96,60)(98,10)
\put(166,98){\vector(4,-1){15}}\put(54,98){\vector(-4,-1){15}}
\put(166,62){\vector(4,1){15}}\put(54,62){\vector(-4,1){15}}
\put(102,66){$0$}
\end{picture}\end{center}

\exercise{5.6.8}{Two stable modes and one unstable mode}
The eigenvalues from Exercise 5.6.2 are $3,4/5,3/5$.
A mode expands along $(1,0,-3)^T$, while the two-dimensional span
of the other eigenvectors is attracted to zero.
The fastest contraction has the smaller modulus $3/5$, along
$(-3,-3,7)^T$; the greatest repulsion has modulus $3$, along
$(1,0,-3)^T$.
\answer{The origin is a saddle. The initial vector in Exercise 2 has coefficient $c_1=2$, so its
trajectory is repelled. Only initial vectors in the stable plane
($c_1=0$) tend to zero.}

\exercise{5.6.9}{Classify the origin by the two eigenvalue moduli}
For $A=\left[\begin{matrix}\frac{17}{10} & - \frac{3}{10}\\- \frac{6}{5} & \frac{4}{5}\end{matrix}\right]$, the characteristic polynomial factors as
\[\frac{\left(\lambda - 2\right) \left(2 \lambda - 1\right)}{2}=0.\]
The eigenpairs, checked by $Av=\lambda v$, are
\[\lambda=\frac{1}{2},\ v=\left[\begin{matrix}1\\4\end{matrix}\right];\qquad
\lambda=2,\ v=\left[\begin{matrix}-1\\1\end{matrix}\right].\]
A decomposition in these two independent directions evolves by
multiplying their coefficients by the corresponding eigenvalues at
each step. Comparing their moduli with $1$ makes the origin an
\textbf{saddle point}. The attracting direction is $\left[\begin{matrix}1\\4\end{matrix}\right]$ (factor $\frac{1}{2}$), and the repelling direction is $\left[\begin{matrix}-1\\1\end{matrix}\right]$ (factor $2$).

\exercise{5.6.10}{Classify the origin by the two eigenvalue moduli}
For $A=\left[\begin{matrix}\frac{3}{10} & \frac{2}{5}\\- \frac{3}{10} & \frac{11}{10}\end{matrix}\right]$, the characteristic polynomial factors as
\[\frac{\left(2 \lambda - 1\right) \left(10 \lambda - 9\right)}{20}=0.\]
The eigenpairs, checked by $Av=\lambda v$, are
\[\lambda=\frac{1}{2},\ v=\left[\begin{matrix}2\\1\end{matrix}\right];\qquad
\lambda=\frac{9}{10},\ v=\left[\begin{matrix}2\\3\end{matrix}\right].\]
A decomposition in these two independent directions evolves by
multiplying their coefficients by the corresponding eigenvalues at
each step. Comparing their moduli with $1$ makes the origin an
\textbf{attractor}. The greatest attraction is along $\left[\begin{matrix}2\\1\end{matrix}\right]$, for the smaller eigenvalue $\frac{1}{2}$. There is no repelling eigendirection.

\exercise{5.6.11}{Classify the origin by the two eigenvalue moduli}
For $A=\left[\begin{matrix}\frac{2}{5} & \frac{1}{2}\\- \frac{2}{5} & \frac{13}{10}\end{matrix}\right]$, the characteristic polynomial factors as
\[\frac{\left(5 \lambda - 4\right) \left(10 \lambda - 9\right)}{50}=0.\]
The eigenpairs, checked by $Av=\lambda v$, are
\[\lambda=\frac{4}{5},\ v=\left[\begin{matrix}5\\4\end{matrix}\right];\qquad
\lambda=\frac{9}{10},\ v=\left[\begin{matrix}1\\1\end{matrix}\right].\]
A decomposition in these two independent directions evolves by
multiplying their coefficients by the corresponding eigenvalues at
each step. Comparing their moduli with $1$ makes the origin an
\textbf{attractor}. The greatest attraction is along $\left[\begin{matrix}5\\4\end{matrix}\right]$, for the smaller eigenvalue $\frac{4}{5}$. There is no repelling eigendirection.

\exercise{5.6.12}{Classify the origin by the two eigenvalue moduli}
For $A=\left[\begin{matrix}\frac{1}{2} & \frac{3}{5}\\- \frac{3}{10} & \frac{7}{5}\end{matrix}\right]$, the characteristic polynomial factors as
\[\frac{\left(5 \lambda - 4\right) \left(10 \lambda - 11\right)}{50}=0.\]
The eigenpairs, checked by $Av=\lambda v$, are
\[\lambda=\frac{4}{5},\ v=\left[\begin{matrix}2\\1\end{matrix}\right];\qquad
\lambda=\frac{11}{10},\ v=\left[\begin{matrix}1\\1\end{matrix}\right].\]
A decomposition in these two independent directions evolves by
multiplying their coefficients by the corresponding eigenvalues at
each step. Comparing their moduli with $1$ makes the origin an
\textbf{saddle point}. The attracting direction is $\left[\begin{matrix}2\\1\end{matrix}\right]$ (factor $\frac{4}{5}$), and the repelling direction is $\left[\begin{matrix}1\\1\end{matrix}\right]$ (factor $\frac{11}{10}$).

\exercise{5.6.13}{Classify the origin by the two eigenvalue moduli}
For $A=\left[\begin{matrix}\frac{4}{5} & \frac{3}{10}\\- \frac{2}{5} & \frac{3}{2}\end{matrix}\right]$, the characteristic polynomial factors as
\[\frac{\left(5 \lambda - 6\right) \left(10 \lambda - 11\right)}{50}=0.\]
The eigenpairs, checked by $Av=\lambda v$, are
\[\lambda=\frac{11}{10},\ v=\left[\begin{matrix}1\\1\end{matrix}\right];\qquad
\lambda=\frac{6}{5},\ v=\left[\begin{matrix}3\\4\end{matrix}\right].\]
A decomposition in these two independent directions evolves by
multiplying their coefficients by the corresponding eigenvalues at
each step. Comparing their moduli with $1$ makes the origin an
\textbf{repeller}. The greatest repulsion is along $\left[\begin{matrix}3\\4\end{matrix}\right]$, for the larger eigenvalue $\frac{6}{5}$. There is no attracting eigendirection.

\exercise{5.6.14}{Classify the origin by the two eigenvalue moduli}
For $A=\left[\begin{matrix}\frac{17}{10} & \frac{3}{5}\\- \frac{2}{5} & \frac{7}{10}\end{matrix}\right]$, the characteristic polynomial factors as
\[\frac{\left(10 \lambda - 13\right) \left(10 \lambda - 11\right)}{100}=0.\]
The eigenpairs, checked by $Av=\lambda v$, are
\[\lambda=\frac{11}{10},\ v=\left[\begin{matrix}-1\\1\end{matrix}\right];\qquad
\lambda=\frac{13}{10},\ v=\left[\begin{matrix}-3\\2\end{matrix}\right].\]
A decomposition in these two independent directions evolves by
multiplying their coefficients by the corresponding eigenvalues at
each step. Comparing their moduli with $1$ makes the origin an
\textbf{repeller}. The greatest repulsion is along $\left[\begin{matrix}-3\\2\end{matrix}\right]$, for the larger eigenvalue $\frac{13}{10}$. There is no attracting eigendirection.

\exercise{5.6.15}{A stochastic model with a nonzero limiting vector}
The columns of $A$ sum to $1$. The given $v_1=(.1,.6,.3)^T$
has $Av_1=v_1$. Eigenvectors for $.5,.2$ are
$v_2=(2,-3,1)^T$ and $v_3=(-1,0,1)^T$.
Solving $c_1v_1+c_2v_2+c_3v_3=(0,.3,.7)^T$ gives
$c_1=1,c_2=1/10,c_3=3/10$. Thus
\[x_k=v_1+\tfrac1{10}(.5)^kv_2+\tfrac3{10}(.2)^kv_3.\]
At zero the coordinates are $(.1+.2-.3,.6-.3,.3+.1+.3)$,
as required. The last two modes tend to zero, so $x_k\to v_1$.
The limiting proportions are $10\%,60\%,30\%$.

\exercise{5.6.16}{The full Hertz three-location solution}
In City Airport, Downtown, Metro Airport order, the matrix from
Section 4.9 is
\[A=\frac1{100}\begin{bmatrix}90&1&9\\1&90&1\\9&9&90\end{bmatrix}.\]
Exact calculation gives three eigenpairs:
\[1,\ (91,19,99)^T;\qquad .89,\ (-1,1,0)^T;
\qquad .81,\ (-1,0,1)^T.\]
The eigenvalues are distinct, so their eigenvectors form a basis.
Therefore the general solution is
\[x_k=c_1(91,19,99)^T+c_2(.89)^k(-1,1,0)^T
+c_3(.81)^k(-1,0,1)^T.\]
The constants are determined by the initial fleet distribution.
For total fleet size $N$, adding the coordinates gives $N=209c_1$,
so the limit is $(N/209)(91,19,99)^T$. At $N=2000$ the downtown
limit is $2000(19)/209\approx181.818$ cars. This recovers the
steady-state interpretation while providing the requested general dynamics.

\exercise{5.6.17}{A two-stage growing population with all four requested graphs}
\step{(a) Translate births and survival into matrix entries.}
Juveniles next year are born to this year's adults, while next year's
adults are surviving juveniles plus surviving adults. Hence
\[A=\begin{bmatrix}0&1.6\\.3&.8\end{bmatrix}.\]
\step{(b) Find growth and the eventual ratio.}
The characteristic polynomial is $\lambda^2-.8\lambda-.48$,
with roots $1.2,-.4$. A dominant eigenvector is $(4,3)^T$;
a second is $(-4,1)^T$. The growing mode therefore has factor
$1.2$, a $20\%$ annual growth rate, and limiting juvenile/adult
ratio $4/3$. For a nonzero nonnegative initial population the growing
coefficient is positive: solving $x_0=c_1(4,3)^T+c_2(-4,1)^T$
gives $c_1=(j_0+4a_0)/16>0$.
\step{(c) Calculate the eight-year trajectory and plot it.}
Beginning at $(15,10)^T$, the actual iterates are
\[\begin{array}{c|rrr}
k&j_k&a_k&j_k/a_k\\\hline
0&15.000000&10.000000&1.500000\\
1&16.000000&12.500000&1.280000\\
2&20.000000&14.800000&1.351351\\
3&23.680000&17.840000&1.327354\\
4&28.544000&21.376000&1.335329\\
5&34.201600&25.664000&1.332668\\
6&41.062400&30.791680&1.333555\\
7&49.266688&36.952064&1.333259\\
8&59.123302&44.341658&1.333358
\end{array}.\]
The four graphs use these points; line segments simply guide the eye
between successive integer years.
\begin{minipage}{.48\linewidth}\centering
\textbf{Juveniles}\par
\setlength{\unitlength}{1pt}\begin{picture}(175,120)
\put(25,18){\vector(1,0){140}}\put(25,18){\vector(0,1){85}}
\put(160,4){$k$}\put(25,6){\makebox(0,0){\scriptsize0}}
\put(153,6){\makebox(0,0){\scriptsize8}}
\put(19,18){\makebox(0,0)[r]{\scriptsize0.00}}
\put(19,93){\makebox(0,0)[r]{\scriptsize63.85}}
\polyline(25.000,35.619)(41.000,36.793)(57.000,41.491)(73.000,45.814)(89.000,51.527)(105.000,58.172)(121.000,66.231)(137.000,75.867)(153.000,87.444)
\put(25.000,35.619){\circle*{2}}\put(41.000,36.793){\circle*{2}}\put(57.000,41.491){\circle*{2}}\put(73.000,45.814){\circle*{2}}\put(89.000,51.527){\circle*{2}}\put(105.000,58.172){\circle*{2}}\put(121.000,66.231){\circle*{2}}\put(137.000,75.867){\circle*{2}}\put(153.000,87.444){\circle*{2}}
\end{picture}\end{minipage}\begin{minipage}{.48\linewidth}\centering
\textbf{Adults}\par
\setlength{\unitlength}{1pt}\begin{picture}(175,120)
\put(25,18){\vector(1,0){140}}\put(25,18){\vector(0,1){85}}
\put(160,4){$k$}\put(25,6){\makebox(0,0){\scriptsize0}}
\put(153,6){\makebox(0,0){\scriptsize8}}
\put(19,18){\makebox(0,0)[r]{\scriptsize0.00}}
\put(19,93){\makebox(0,0)[r]{\scriptsize47.89}}
\polyline(25.000,33.661)(41.000,37.577)(57.000,41.179)(73.000,45.940)(89.000,51.477)(105.000,58.193)(121.000,66.224)(137.000,75.871)(153.000,87.444)
\put(25.000,33.661){\circle*{2}}\put(41.000,37.577){\circle*{2}}\put(57.000,41.179){\circle*{2}}\put(73.000,45.940){\circle*{2}}\put(89.000,51.477){\circle*{2}}\put(105.000,58.193){\circle*{2}}\put(121.000,66.224){\circle*{2}}\put(137.000,75.871){\circle*{2}}\put(153.000,87.444){\circle*{2}}
\end{picture}\end{minipage}\par\begin{minipage}{.48\linewidth}\centering
\textbf{Total population}\par
\setlength{\unitlength}{1pt}\begin{picture}(175,120)
\put(25,18){\vector(1,0){140}}\put(25,18){\vector(0,1){85}}
\put(160,4){$k$}\put(25,6){\makebox(0,0){\scriptsize0}}
\put(153,6){\makebox(0,0){\scriptsize8}}
\put(19,18){\makebox(0,0)[r]{\scriptsize0.00}}
\put(19,93){\makebox(0,0)[r]{\scriptsize111.74}}
\polyline(25.000,34.780)(41.000,37.129)(57.000,41.357)(73.000,45.868)(89.000,51.506)(105.000,58.181)(121.000,66.228)(137.000,75.869)(153.000,87.444)
\put(25.000,34.780){\circle*{2}}\put(41.000,37.129){\circle*{2}}\put(57.000,41.357){\circle*{2}}\put(73.000,45.868){\circle*{2}}\put(89.000,51.506){\circle*{2}}\put(105.000,58.181){\circle*{2}}\put(121.000,66.228){\circle*{2}}\put(137.000,75.869){\circle*{2}}\put(153.000,87.444){\circle*{2}}
\end{picture}\end{minipage}\begin{minipage}{.48\linewidth}\centering
\textbf{Juvenile/adult ratio}\par
\setlength{\unitlength}{1pt}\begin{picture}(175,120)
\put(25,18){\vector(1,0){140}}\put(25,18){\vector(0,1){85}}
\put(160,4){$k$}\put(25,6){\makebox(0,0){\scriptsize0}}
\put(153,6){\makebox(0,0){\scriptsize8}}
\put(19,18){\makebox(0,0)[r]{\scriptsize1.26}}
\put(19,93){\makebox(0,0)[r]{\scriptsize1.52}}
\polyline(25.000,87.231)(41.000,23.769)(57.000,44.351)(73.000,37.429)(89.000,39.730)(105.000,38.962)(121.000,39.218)(137.000,39.133)(153.000,39.161)
\put(25.000,87.231){\circle*{2}}\put(41.000,23.769){\circle*{2}}\put(57.000,44.351){\circle*{2}}\put(73.000,37.429){\circle*{2}}\put(89.000,39.730){\circle*{2}}\put(105.000,38.962){\circle*{2}}\put(121.000,39.218){\circle*{2}}\put(137.000,39.133){\circle*{2}}\put(153.000,39.161){\circle*{2}}
\end{picture}\end{minipage}
The ratio oscillates toward $4/3$ because the secondary eigenvalue
is negative. It is already close by years 3--4 and differs by less
than $.001$ from $4/3$ by year 5. The underlying relative mode
shrinks by $|-.4/1.2|=1/3$ each year.
A complete plotting program, including the requested total and ratio
panels, is:
\begin{verbatim}
import numpy as np
import matplotlib.pyplot as plt
A = np.array([[0.,1.6],[.3,.8]])
x = np.array([15.,10.])
values = [x.copy()]
for k in range(8):
    x = A @ x
    values.append(x.copy())
v = np.array(values)
series = [v[:,0],v[:,1],v.sum(axis=1),v[:,0]/v[:,1]]
labels = ['Juveniles','Adults','Total','Juveniles/adults']
fig, axes = plt.subplots(2,2)
for ax, y, label in zip(axes.flat, series, labels):
    ax.plot(range(9), y, 'o-')
    ax.set(xlabel='Year', ylabel=label)
fig.tight_layout()
plt.show()
\end{verbatim}

\exercise{5.6.18}{A three-stage buffalo model}
\step{(a) Construct the stage matrix.}
Only adults give birth. Calves survive into the yearling stage and
yearlings into adulthood. Thus
\[A=\begin{bmatrix}0&0&.42\\.60&0&0\\0&.75&.95\end{bmatrix}.\]
\step{(b) Compute the dominant mode and its composition.}
The characteristic polynomial is
$\lambda^3-.95\lambda^2-.189$. Its real positive root is
$\lambda_1\approx1.104834392$; the other two roots are complex
conjugates of modulus less than $1$. Since $\lambda_1>1$, the
nonzero nonnegative population grows asymptotically by factor
$1.104834392$, about $10.483439\%$ per year.
Normalize the dominant eigenvector to 100 adults. The first row of
$Av=\lambda_1v$ gives $c=42/\lambda_1$, and the second gives
$y=.6c/\lambda_1$. Therefore the eventual composition per 100
adults is approximately
\[(c,y,a)=(38.014747,20.644586,100).\]
Substituting these formulas into the third row reproduces the
characteristic equation; this checks the eigenvector normalization.
\step{Show that every nonzero nonnegative initial population has a growing mode.}
An initial state could in principle have zero coefficient in a particular
eigenvector direction, so we must verify that the dominant coefficient
is positive here. A useful row multiplier is the positive left eigenvector
\[
w=\begin{bmatrix}.45/\lambda_1^2\\.75/\lambda_1\\1\end{bmatrix}.
\]
Indeed direct multiplication gives
\[
w^TA=\begin{bmatrix}.45/\lambda_1&.75&.189/\lambda_1^2+.95\end{bmatrix}
=\lambda_1w^T.
\]
The last equality uses the characteristic equation, divided by
$\lambda_1^2$: $\lambda_1=.95+.189/\lambda_1^2$.
All three entries of $w$ are strictly positive because $\lambda_1>0$.

The three distinct eigenvalues give an eigenvector expansion
$x_0=c_1v_1+c_2v_2+c_3v_3$, allowing complex coefficients for the
conjugate pair. For $j=2,3$, calculate the same scalar in two ways:
\[
w^TAv_j=\lambda_1w^Tv_j=\lambda_jw^Tv_j.
\]
Since $\lambda_j\ne\lambda_1$, this forces $w^Tv_j=0$.
Multiplying the expansion of $x_0$ by $w^T$ therefore isolates
\[
c_1=\frac{w^Tx_0}{w^Tv_1}>0.
\]
The numerator is positive whenever $x_0$ has nonnegative entries and at
least one positive entry; the denominator is positive because $v_1$
has the positive components computed above. Consequently
\[
\frac{x_k}{\lambda_1^k}
=c_1v_1+c_2\left(\frac{\lambda_2}{\lambda_1}\right)^kv_2
       +c_3\left(\frac{\lambda_3}{\lambda_1}\right)^kv_3
\longrightarrow c_1v_1.
\]
Both remaining ratios have modulus less than $1$, so their powers tend
to zero. Thus the dominant mode really is present for every nonzero
nonnegative initial population, justifying both the growth rate and the
eventual stage proportions claimed above.
